\documentclass[11pt]{article}
\usepackage[a4paper,margin=1in]{geometry}
\usepackage{amsmath,amssymb,amsthm,hyperref}
\newtheorem{theorem}{Theorem}
\newtheorem{lemma}{Lemma}
\newcommand{\Mat}{\operatorname{Mat}}
\newcommand{\rank}{\operatorname{rank}}
\newcommand{\GL}{\operatorname{GL}}
\title{A uniform point-mass bound for sums of prescribed-rank matrices}
\author{Principia Math In Progress}
\date{22 September 2026}
\begin{document}
\maketitle
\begin{abstract}
Independent uniform matrices of prescribed ranks over a finite field
have a uniformly bounded density relative to the uniform measure as soon
as the sum of the ranks is at least twice the matrix dimension.
An explicit absolute bound proves Conjecture 5.5 of Gill, Lifshitz,
Pyber and Szab\'o, with $K(\epsilon)=\lceil2/\epsilon\rceil$.
The proof uses rectangular factorization and a rank-sensitive estimate
for the probability of full rank.
\end{abstract}

\section*{Task 1: the complete result}
Let $q$ be any prime power, $n\ge1$, and let $R_{n,q}(r)$ denote the
number of rank-$r$ matrices in $\Mat_{n\times n}(\mathbb F_q)$.
For a fixed matrix $M$ of rank $t$, let
$N_{n,q}(r_1,\ldots,r_k;t)$ count tuples $(X_1,\ldots,X_k)$ with
$\rank X_i=r_i$ and $\sum_iX_i=M$.
This number depends only on $t$, by invertible left and right multiplication.

\begin{theorem}\label{main}
For every prime power $q$, every $n\ge1$, and every sequence of integers
$0\le r_i\le n$ such that $\sum_{i=1}^k r_i\ge2n$, one has
\[
 N_{n,q}(r_1,\ldots,r_k;t)
 \le 32768\,q^{-n^2}\prod_{i=1}^k R_{n,q}(r_i)
 \qquad(0\le t\le n).
\]
Consequently Conjecture 5.5, equivalently Question 5.22, of
\cite{source} holds with $K(\epsilon)=\lceil2/\epsilon\rceil$ and an
absolute implied constant independent even of $\epsilon$.
\end{theorem}

The conclusion is the upper bound requested in the source. It does not
assert a two-sided estimate or convergence to the uniform distribution.
For example, when $n=1$, $q=2$ and every $r_i=1$, the sum is deterministic.
This example is entirely consistent with the asserted upper bound.

\section*{Two elementary estimates}
The full-rank and subspace-count formulas below are standard and are
also recorded in Fact 5.23 of \cite{source}. We include their estimates
to make the argument self-contained.

\begin{lemma}\label{conditioning}
If $U$ is uniform in $\Mat_{n\times r}(\mathbb F_q)$, where $0\le r\le n$,
then
\[
 p_n(r):=\mathbb P(\rank U=r)
 =\prod_{i=0}^{r-1}(1-q^{i-n})\ge4^{-r/n}.
\]
The same probability applies to full row rank in $\Mat_{r\times n}$.
\end{lemma}
\begin{proof}
Choose the columns successively. After $i$ independent columns their
span has $q^i$ elements, proving the product formula.
The numbers $\log(1-q^{i-n})$, $0\le i<n$, are decreasing in $i$.
Their first $r$ terms therefore have average at least their overall average:
\[
 \frac{\log p_n(r)}r\ge\frac{\log p_n(n)}n\qquad(r\ge1).
\]
Moreover, using $\prod_j(1-a_j)\ge1-\sum_j a_j$ for $0\le a_j\le1$,
\[
 p_n(n)=\prod_{j=1}^n(1-q^{-j})
 \ge\frac12\prod_{j=2}^n(1-2^{-j})\ge\frac14.
\]
The empty product covers $n=1$ in the last estimate and $r=0$ in the
statement. Exponentiating proves the claim; transposition gives its row version.
\end{proof}

\begin{lemma}\label{product}
Let $R\ge2n$. If $A\in\Mat_{n\times R}(\mathbb F_q)$ and
$B\in\Mat_{R\times n}(\mathbb F_q)$ are independent and uniform, then
\[
 \mathbb P(AB=M)\le8q^{-n^2}
 \quad\text{for every }M\in\Mat_{n\times n}(\mathbb F_q).
\]
\end{lemma}
\begin{proof}
Conditional on $A$ having rank $n-d$, the columns of $AB$ are
independent uniform elements of $\operatorname{im}A$.
Thus their probability of equaling the prescribed columns of $M$ is
either zero or $q^{-n(n-d)}$.

There are at most $4q^{d(n-d)}$ subspaces of dimension $d$ in
$\mathbb F_q^n$. Indeed, count ordered independent $d$-tuples by at most
$q^{nd}$, and divide by their number in $\mathbb F_q^d$, which is
$q^{d^2}p_d(d)\ge q^{d^2}/4$ for $d\ge1$.
The same upper bound is valid for $d=0$.
For any fixed such subspace $W$, the probability that $W\subseteq\ker A^{\mathsf T}$
is $q^{-dR}$: each of the $R$ independent columns must satisfy $d$
independent linear equations. A union bound therefore gives
\[
 \mathbb P(\rank A=n-d)
 \le4q^{-d(R-n+d)}\qquad(0\le d\le n).
\]
It follows that
\begin{align*}
 \mathbb P(AB=M)
 &\le\sum_{d=0}^n4q^{-d(R-n+d)}q^{-n(n-d)}\\
 &=4q^{-n^2}\sum_{d=0}^n q^{-d(R-2n+d)}\\
 &\le4q^{-n^2}\sum_{d=0}^{\infty}2^{-d^2}
 \le8q^{-n^2},
\end{align*}
where $R\ge2n$, $q\ge2$ and $d^2\ge d$ were used.
\end{proof}

\section*{Proof of the uniform prescribed-rank bound}
\begin{proof}[Proof of Theorem \ref{main}]
Choose $m$ minimally so that $R=\sum_{i=1}^m r_i\ge2n$.
Then $2n\le R<3n$, since the preceding sum is less than $2n$ and
$r_m\le n$.
Independently choose uniform matrices
\[
 U_i\in\Mat_{n\times r_i}(\mathbb F_q),\qquad
 V_i\in\Mat_{r_i\times n}(\mathbb F_q),\qquad 1\le i\le m,
\]
and let $E$ be the event that every $U_i$ has full column rank and
every $V_i$ has full row rank. Lemma \ref{conditioning} gives
\[
 \mathbb P(E)=\prod_{i=1}^m p_n(r_i)^2
 \ge16^{-R/n}\ge16^{-3}.
\]

Conditional on $E$, the products $U_iV_i$ are independent uniform
rank-$r_i$ matrices. To see uniformity, for any fixed rank-$r$ matrix $X$,
the columns of $U$ must be an ordered basis of $\operatorname{im}X$.
There are exactly $|\GL(r,q)|$ choices, and each determines a unique
full-row-rank $V$ with $UV=X$. The rank-zero case has one empty
factorization. The conditioning event separates across the independent
blocks, so independence also persists.

Before conditioning, form the matrices
\[
 A=[\,U_1\ \cdots\ U_m\,],\qquad
 B=\begin{bmatrix}V_1\\ \vdots\\ V_m\end{bmatrix}.
\]
They are independent uniform matrices of the sizes in Lemma \ref{product},
and $AB=\sum_{i=1}^m U_iV_i$. Consequently, for independent uniform
rank-$r_i$ matrices $X_i$,
\[
 \mathbb P\left(\sum_{i=1}^mX_i=M\right)
 =\mathbb P(AB=M\mid E)
 \le\frac{\mathbb P(AB=M)}{\mathbb P(E)}
 \le8\cdot16^3q^{-n^2}=32768q^{-n^2}.
\]
Adding the remaining independent matrices cannot increase the largest
point probability: average this bound over their sum. Finally, multiply
by $\prod_iR_{n,q}(r_i)$ to obtain the stated count.

If $r_i\ge\epsilon n$ and $k\ge\lceil2/\epsilon\rceil$, then
$\sum_i r_i\ge k\epsilon n\ge2n$. This proves the conjecture for
$0<\epsilon\le1$; for $\epsilon>1$ its hypotheses are impossible.
\end{proof}

The rank-sensitive conditioning bound is what keeps the constant absolute.
Using only $p_n(r)\ge1/4$ separately for each block would introduce an
unnecessary factor exponential in the number of summands.
The source's established case of three ranks at least $2n/3$
is credited to its Theorem 5.3. The present argument proves the full
many-rank extension and the stronger sum-of-ranks criterion.

\section*{Proposed follow-up Tasks 2--5}
These are proposed tasks, not additional completed results.
\begin{enumerate}
\setcounter{enumi}{1}
\item Determine the smallest absolute constant valid under the sole
condition $\sum_i r_i\ge2n$.
\item Find sharp necessary rank-sum thresholds for uniform point-mass
bounds, including sequences with $\sum_i r_i<2n$.
\item Prove the corresponding optimal theorem for rectangular
$a\times b$ matrices with arbitrary prescribed ranks.
\item Characterize the rank sequences and fields for which the normalized
point probabilities converge uniformly to one, and obtain quantitative rates.
\end{enumerate}
\begin{thebibliography}{9}
\bibitem{source} N. Gill, N. Lifshitz, L. Pyber and E. Szab\'o,
\emph{Initiating the proof of the Liebeck--Nikolov--Shalev conjecture},
\href{https://arxiv.org/abs/2408.07800v2}{arXiv:2408.07800v2}, 24 September 2024.
Conjecture 5.5, Question 5.22, Theorem 5.3 and Fact 5.23.
\end{thebibliography}
\end{document}
